\documentclass[10pt,a4paper]{amsart}
\usepackage[margin=19mm]{geometry}
\usepackage{amsmath,amssymb,mathtools}
\usepackage[scr=rsfs,cal=euler]{mathalfa}
\usepackage{xfrac}
\usepackage[unicode,hidelinks,bookmarks=false]{hyperref}
\newcommand{\ie}{i.e.\ }
\newcommand{\cf}{cf.\ }
\newcommand{\resp}{respectively\ }
\newcommand{\Subsection}{\subsection}
\providecommand{\nobreakdash}{\nobreak\hbox{-}}
\setlength{\emergencystretch}{2em}
\multlinegap0pt
\usepackage{mathrsfs}
\allowdisplaybreaks
\newtheorem{theorem}{Theorem}[section]
\newtheorem{lemma}[theorem]{Lemma}
\newtheorem{proposition}[theorem]{Proposition}
\newtheorem{definition}[theorem]{Definition}
\newtheorem{example}[theorem]{Example}
\newtheorem{corollary}[theorem]{Corollary}
\newtheorem{remark}[theorem]{Remark}
\begin{document}
\title[Functions and Means of Accretive Operators]{Functions and Means of Accretive Operators}
\author{47A64; Secondary 15A15 ạte 6/29, 2026; Mitsuru Uchiyama}
\date{}
\maketitle
\noindent\emph{Source and licence.} Functions and Means of Accretive Operators. Version arXiv:2607.09152v1 (CC BY 4.0). This material is distributed under the Creative Commons Attribution 4.0 International licence, http://creativecommons.org/licenses/by/4.0/, as recorded in the arXiv OAI licence metadata for this exact version and shown on the abstract page https://arxiv.org/abs/2607.09152v1. Full rights notice: licenses/arxiv-2607.09152-CC-BY-4.0.txt.\par\smallskip
\noindent\textit{Editorial scope.} Counted unit: the original \texttt{theorem} environment \texttt{Th:3-1} at source lines 417--429 together with its complete proof at lines 430--468; the delivered document keeps source lines 115--477 so that every referenced label and every citation resolves inside it. Native editable source is the paper's own concatenated TeX; the AMS wrapper is editorial formatting only. No publisher class, upstream script or unrelated artwork is redistributed.
\begin{equation}
\begin{split}
&f(t)= f(0) + bt + \int_0^{\infty} (\frac{1}{\lambda} - \frac{1}{\lambda +t}) d\mu(\lambda), \\ 
&b\geqq 0, \quad \int_0^{\infty}\frac{1}{\lambda^2 + 1} d\mu (\lambda) <\infty, \quad \int_0^1 
\frac{1}{\lambda} d\mu(\lambda)<\infty. \label{eq:1}
\end{split}
\end{equation}
For example,  for $0<r<1$ and $0<t<\infty$
$$t^r = \frac{\sin r \pi}{\pi} \int_0^{\infty} \frac{t}{\lambda + t } \lambda ^{r-1} d\lambda =
 \frac{\sin r \pi}{\pi} \int_0^{\infty} 
 (\frac{1}{\lambda} - \frac{1}{\lambda +t}) \lambda ^{r} d\lambda.$$
(refer to \cite{D}, \cite{B-M}, \cite{Simon} for more details)\par
Let $T$ be an operator with ${\it Sp}(T) \cap (-\infty, 0] = \emptyset$ and 
 $\gamma$ a smooth curve 
in $\mathbb{C} \smallsetminus (-\infty, 0]$ and surrounding 
$\it{Sp}(T)$. Then 
the Riesz-Dunford integral
$$f(T) = \frac{1}{2\pi i}\int_{\gamma}  f(z) (z I - T)^{-1} dz$$
 is well-defined for an analytic extension $f(z)$ of $f(t)$ given \eqref{eq:1}
(see Chapter VII of \cite{D-S}).
We will show that
$$ f(T)=  f(0) I  + b T + \int_0^{\infty} \left(\frac{1}{\lambda } I - (\lambda I  + T)^{-1}\right)d\mu(\lambda), $$
and that this holds for an accretive operator $A$ even if $0\in {\it Sp}(A)$.
This naturally deduces 
$$A^r = \frac{\sin r \pi}{\pi} \int_0^{\infty} 
 (\frac{1}{\lambda} I - (\lambda I + A)^{-1}) \lambda ^{r} d\lambda \quad (0<r<1),$$
which is known as the Balakrishnan formula\cite{Bal} for a closed and accretive operator
 (cf. \cite{N-F}, \cite{Kato-book}).
  So, the formula $f(T)$ mentioned above is a generalization of fractional powers. \par
 Drury \cite{Drury}
defined {\it geometric mean} $A\# B$ of strictly accretive {\it matrices} $A$ and $B$ by 
\begin{equation}\label{eq:1-0}
(A\# B)^{-1}:= \frac{2}{\pi} \int_0^{\infty} (tA + t^{-1} B)^{-1} t^{-1} d t.
\end{equation}
and showed $A \# B = A^{1/2} (A^{-1/2} B A^{-1/2})^{1/2} A^{1/2}$. We extend this as follows: 
Let $A$ and $B$ be accretive operators such that  
${\it Sp}(BA^{-1}) \cap (-\infty, 0] = \emptyset$.  
Then we define the geometric mean and the general operator mean $\sigma_f$ by 
$$A\#B:=A^{1/2} (A^{-1/2} B A^{-1/2})^{1/2} A^{1/2}, \quad A\sigma_f B:= A^{1/2} 
f( A^{-1/2} B A^{-1/2}) A^{1/2},$$
and show 
\begin{equation}\label{eq:1-1}
A\#B= \frac{1}{\pi}\int_0^{\infty} \frac{1}{\lambda} 
A ( \lambda A + B)^{-1} B \sqrt{\lambda} d\lambda.\\
%A\sigma_f B=  f(0) A  + b B + \int_0^{\infty} 
%\frac{1}{\lambda } A ( \lambda A + B)^{-1} B d\mu(\lambda).\notag
\end{equation}
A new formula  
$$A (A\# B)^{-1} B= A \# B$$
leads us to see that  \eqref{eq:1-1}  is coincident with \eqref{eq:1-0}. 
One of the objective of this paper is to show 
$$A\#B = A^{1/2} B^{1/2}$$
if $A$ and  $B$ are normal and commutative. Another one is to verify  
$$0\leqq A^*\nabla A \leqq A^* \# A \leqq A^* \,!\, A,$$
where symbols $\nabla, !$ conventionally stand for arithmetic and harmonic mean, 
respectively. At first sight this seems strange, because 
\begin{center}
$A \nabla B \geqq A \# B \geqq A \,!\, B\geqq 0$ for $A, \,B \geqq 0$.
\end{center}
However we should remember that $0\leqq {\it Re}(\alpha) \leqq |\alpha| \leqq 
|\alpha |^2 /{\it Re}(\alpha)$ for a complex number $\alpha$ with 
${\it Re}(\alpha) >0$.
We further show that $A^* \# A = | A |$ if and only if $A$ is normal. Consequently, for 
a stricly accretive and normal operator $A$  
$$|A|=  \frac{1}{\pi}\int_0^{\infty}A (\lambda A + A^*)^{-1} A^* 
\lambda^{-1/2} d \lambda , $$
which deduces the simplest case:
$$1 = \frac{1}{\pi} \int_0^{\infty} \frac{1}{(\lambda e^{i\theta}  +  e^{-i\theta})
\sqrt{\lambda} } d \lambda \quad (|\theta | < \pi/2). $$
This is a basic formula shown by using residue.
For an accretive and normal operator $A$ we also gain 
$$A + A^* \leqq   A^{1-r} A^{*r} + A^r A^{*(1-r)} \quad (0\leqq r \leqq 1).$$

\section{Preliminaries}
Let $f(t)$ be an operator monotone function on $(0, \infty)$ represented by 
$\eqref{eq:1}$. Then by Loewner's theorem 
\begin{equation}
f(z) = f(0) + b z + \int_0^{\infty} (\frac{1}{\lambda} - \frac{1}{\lambda + z}) d\mu(\lambda) \label{eq:2}
\end{equation}
 is the analytic extension of $f(t)$ to $\mathbb{C} \smallsetminus (-\infty, 0]$. 
But it is possible to see, without using Loewner's theorem,  that $f(t)$ is of the 
form $\eqref{eq:1}$ if and only if 
$f(t)$ is operator monotone on $(0, \infty)$ with $f(0)>-\infty$; for instance see \cite{Hansen}. 
Just in case, we give a proof of it.
$\arg z$ always stands for the principal argument of $z\ne 0$, namely 
$-\pi < \arg z \leqq \pi$.
\begin{lemma}\label{Lem:2-1} Let $f(t)$ be an operator monotone function on $(0, \infty)$ given 
by \eqref{eq:1}. Then 
$f(z)$ given by \eqref{eq:2} is the analytic extension of $f(t)$ to 
$\mathbb{C}\setminus (-\infty, 0]$ such that $\overline{f(\overline{z})} = f(z)$.
 Further, 
for any $\theta: \frac{\pi}{2} \leqq \theta < \pi$, 
$f(z)$ is continuous 
on $$\{ z \in \mathbb{C}: z=0\,\, \text{or}\,\, |\arg z| \leqq \theta \}.$$
\end{lemma}
\begin{proof} The integrand in \eqref{eq:2} is continuous on  $\{(\lambda , z ): 0<\lambda <\infty, 
z\in \mathbb{C} \smallsetminus (-\infty, 0] \}$, 
 and integrable with respect to $\mu(\lambda)$; indeed, 
$$\lim_{\lambda \to 0} \lambda |\frac{1}{\lambda} - \frac{1}{\lambda -z}| =1, \quad 
\lim_{\lambda \to \infty} ( 1+ \lambda ^2) |\frac{1}{\lambda} 
- \frac{1}{\lambda -z}| =|z|.$$
To see the analyticity, take $z_0 \in \mathbb{C} \smallsetminus (-\infty, 0] $ and put
$0<\delta < {\it dist}\{ z_0, (-\infty, 0]\}$. For $z$ s.t. $|z - z_0|<\delta$, 
$$\frac{f(z) - f(z_0)}{z- z_0} = b + \int_0^{\infty} (\frac{1}{(\lambda + z_0) (\lambda + z)} d\mu(\lambda).$$
Since 
\begin{center} 
$|\frac{1}{(\lambda + z_0) (\lambda + z)}| \leqq \frac{1}{|\lambda + z_0| (|\lambda + z_0| - \delta )} $
\end{center} 
and there is $M>0$ such that 
$\frac{1}{|\lambda + z_0| (|\lambda + z_0| - \delta )} \leqq M \frac{1}{1 + \lambda ^2}$ for $\lambda >0$,
we therefore get 
$$ f'(z)= b + \int_0^{\infty} (\frac{1}{ (\lambda + z)^2} d\mu(\lambda).$$
 We next show the continuity of $f(z)$ at $z=0$.
Let $z\ne 0$, $|\arg z| \leqq \theta$ and $|z|\leqq 1$. Then
\begin{center}
$|\frac{z}{\lambda + z}|\leqq 1\, \,\, (\Re{z} \geqq 0), \quad 
|\frac{z}{\lambda + z}|\leqq |\frac{z}{\Im z}|\leqq \frac{1}{\sin \theta}\,\,
 \, (\Re{z} <0)$.
\end{center}
Consequently  
$$|\frac{1}{\lambda} - \frac{1}{\lambda + z} | \leqq \frac{1}{\lambda \sin \theta} \quad (\lambda > 0).$$
On the other hand, we have
$$|\frac{1}{\lambda} - \frac{1}{\lambda + z} | \leqq \frac{2}{\lambda ^2} \quad 
(\lambda \geqq 2).$$ Combining these inequalities lead us to   
$$|\frac{1}{\lambda} - \frac{1}{\lambda + z} | \leqq g(\lambda)
: =
\begin{cases}
\frac{1}{\lambda} \frac{1}{ |\sin \theta|}  &  (0 <\lambda \leqq 2)\\
\frac{2}{\lambda^2} \frac{1}{ |\sin \theta|}  & ( 2\leqq \lambda <\infty).
\end{cases}
$$
Since $g(\lambda)$ is continuous and integrable,
$$\lim_{z\to 0} f(z) = f(0) + \int_0^{\infty} \lim_{z\to 0} (\frac{1}{\lambda} - 
\frac{1}{\lambda + z})
 d\mu(\lambda) = f(0).$$
\end{proof}

Let $T$ be an operator such that ${\it Sp}(T) \cap (-\infty, 0] = \emptyset$ and 
$f(z)$ an analytic function expressed by \eqref{eq:2}. Then  
$f(T)$ is defined by the Riesz-Dunford integral. 
Recall that ${\it Sp}(f(T)) = f ({\it Sp}(T))$ and $f(g(T))= (f\circ g)(T)$, so 
${\it Sp}(T^{r}) \subseteqq \{z: -r\pi <\arg z <r \pi\} $,  and  $(T^{r})^s = T^{rs}$,    
where $0<r, s <1$, and  
$ (T^{1/n})^n = T$, furthermore, we have $(T^{-1})^r = (T^r)^{-1}$ for $0<r<1$.  
We here represent $f(T)$ by using \eqref{eq:2}, that is \eqref{eq:1}.  

\begin{proposition}\label{Prop:2-2}
Let $T$ be an operator with ${\it Sp}(T) \cap (-\infty, 0] = \emptyset$, and $f(t)$ an
 operator monotone function on $(0, \infty)$ represented by \eqref{eq:1}, namely 
\eqref{eq:2}. 
Define $f(T)$ by the Riesz-Dunford integral.   
Then, for all $\mathbf{x}, \mathbf{y}$
$$
 (f(T) \mathbf{x}, \mathbf{y})=  f(0) (\mathbf{x}, \mathbf{y})  + b (T \mathbf{x}, \mathbf{y})
 + \int_0^{\infty} \left(\frac{1}{\lambda}(\mathbf{x}, \mathbf{y}) - 
((\lambda I  + T)^{-1}\mathbf{x}, \mathbf{y})\right)  d\mu(\lambda).$$
 We write this as
\begin{equation}
 f(T)=  f(0) I  + b T + \int_0^{\infty} \left(\frac{1}{\lambda } I - (\lambda I  + T)^{-1}\right) d\mu(\lambda).\label{eq:3}
\end{equation}
In this case we have
$$|| f(T)|| \leqq   |f(0)|  + b ||T|| + \int_0^{\infty} ||\frac{1}{\lambda } I - (\lambda I  + T)^{-1}|| d\mu(\lambda)
<\infty.$$
Moreover,  $f(T)^* = f(T^*)$,  and  $f(T)$ is normal if $T$ is.

\end{proposition}
 
\begin{proof}
Let $\gamma$ be a smooth, simple and closed curve in $\mathbb{C} \smallsetminus (-\infty, 0]$ 
 surrounding $\it{Sp}(T)$. Then 
$$f(T) = \frac{1}{2\pi i} \int_{\gamma}  f(z) (z I - T)^{-1} dz.$$
Since $\frac{1}{2\pi i} \int_{\gamma}  (f(0) + b z) (z I - T)^{-1} dz = f(0) I  + b T$, we need to show 
\begin{gather*}
\frac{1}{2\pi i} \int_{\gamma} \left(\int_0^{\infty} (\frac{1}{\lambda } - \frac{1}{\lambda + z}) d\mu(\lambda) 
((z I - T)^{-1} \mathbf{x}, \mathbf{y})\right) dz\\
= \int_0^{\infty} ((\frac{1}{\lambda } I - (\lambda I  + T)^{-1}) \mathbf{x}, \mathbf{y})  d\mu(\lambda).
\end{gather*}
By the preceding lemma, the left side exists. To use the Fubini theorem
 represent $\gamma$ as 
$ z=\phi (s), \alpha \leqq s \leqq \beta$. Then the left side is equal to 
$$\frac{1}{2\pi i} \int_{\alpha}^{\beta} \left(\int_0^{\infty} (\frac{1}{\lambda} - \frac{1}{\lambda + \phi (s)})
d\mu(\lambda) 
((\phi (s) I - T)^{-1} \mathbf{x}, \mathbf{y}) \right) \phi'(s) d s,$$
which is finite. 
The integrand of the double integral 
$$(\frac{1}{\lambda} - \frac{1}{\lambda + \phi (s)}) ((\phi (s) I - T)^{-1} \mathbf{x}, \mathbf{y}) \phi'(s)$$
is continuous with respect to  $\lambda$ and $s$. Hence the above successive integral is equal to 
\begin{align*}
&\int_0^{\infty}d\mu(\lambda) 
\frac{1}{2\pi i} \int_{\alpha}^{\beta} (\frac{1}{\lambda} - \frac{1}{\lambda + \phi (s)}) ((\phi (s) I - T)^{-1}
 \mathbf{x}, \mathbf{y})  \phi'(s) d s\\
= & \int_0^{\infty}d\mu(\lambda) \frac{1}{2\pi i} \int_{\gamma} (\frac{1}{\lambda} 
- \frac{1}{\lambda + z})
 ((z I - T)^{-1} \mathbf{x}, \mathbf{y})  d z\\
=&  \int_0^{\infty}((\frac{1}{\lambda}I -
(\lambda I + T)^{-1}) \mathbf{x}, \mathbf{y}) d \mu(\lambda). 
\end{align*}
We thus got the required formula.  
Since 
$$\lim_{\lambda \to + 0} \lambda ||\frac{1}{\lambda} I - (\lambda I  + T)^{-1}|| = 1\, \, \text{and} \, \,
\lim_{\lambda \to + \infty} \lambda^2 ||\frac{1}{\lambda} I - (\lambda I  + T)^{-1}|| = || T||,$$ 
%$\lim_{\lambda \to + 0} \lambda ||\frac{1}{\lambda} I - (\lambda I  + T)^{-1}|| = 1$
% and 
%$\lim_{\lambda \to + \infty} \lambda^2 ||\frac{1}{\lambda} I - (\lambda I  + T)^{-1}|| = || T||.$, 
 $$\int_0^{\infty} ||\frac{1}{\lambda} I - (\lambda I  + T)^{-1}|| d\mu(\lambda)
<\infty.$$
From \eqref{eq:3} it follows that $f(T)^* = f(T^*)$, and this lead us to the last statement.
\end{proof}

\begin{remark}\label{Re:2-1}\rm
One may think the integral in $\eqref{eq:3}$ as the Bochner integral of the operator valued 
function. Indeed, the integrand is continuous, and 
$$\int_0^{\infty} ||\frac{1}{\lambda} I - (\lambda I  + T)^{-1}|| d\mu(\lambda)
<\infty.$$
Consequently the integral is well-defined as the Bochner integral and 
\begin{align*}
&\left(\int_0^{\infty} (\frac{1}{\lambda} I - (\lambda I  + T)^{-1}) d\mu(\lambda)
\mathbf{x}, \mathbf{y}\right)\\
 =& \int_0^{\infty} \left(\frac{1}{\lambda}(\mathbf{x}, \mathbf{y}) - 
((\lambda I  + T)^{-1}\mathbf{x}, \mathbf{y})\right)  d\mu(\lambda)
\end{align*}
(cf. Section V. 5 of \cite{Y}). 
But it is enough for us to consider $\eqref{eq:3}$ in the weak sense (see 3.26 Definition of \cite{Rudin}). 
This kind of representation was dealt in \cite{Sano}
\end{remark}
\begin{remark}\label{Re:2-2}\rm
(i). Let $T$ be an $n\times n$ normal matrix such that ${\it Sp}(T) \cap (-\infty, 0] = \emptyset$ and $\alpha_1 P_1 + \cdots + \alpha_k P_k  
\, (1\leqq k \leqq n)$ its spectral decomposition. Then  
$$f(T)= \int_{\gamma} f(z) (z-T)^{-1} dz =  f(\alpha_1) P_1 + \cdots + 
f(\alpha_k) P_k.$$
(ii). Let $T$ be a normal operator such that ${\it Sp}(T) \cap (-\infty, 0] = \emptyset$ and 
$\int_{{\it Sp}(T)} z d E_z$ its spectral decomposition. Then the Fubini 
theorem entails that the Riesz-Dunford integral $f(T)$ coincides with the functional calculus 
$\int_{{\it Sp}(T)} f(z) d E_z$.\\
(iii). 
Let $T= \int_{{\it Sp}(T)} z dE_z$ be a normal operator with ${\it Sp}(T) \cap (-\infty, 0) = \emptyset$. Then for $\epsilon>0$, the Riesz-Dunford integral $f(T+ \epsilon I)$ is well-defined. 
By taking $\theta$ so that  
${\it Sp}(T) \subseteq \{ z \in \mathbb{C}: z=0\,\, \text{or}\,\, |\arg z| \leqq \theta \}$, we obtain 
$$||f(T+ \epsilon I) - \int_{{\it Sp}(T)} f(z) dE_z|| \leqq 
 \sup \{|f(z+ \epsilon) - f(z)|: z \in {\it Sp}(T)\} \to 0 \,\, (\epsilon \to 0).$$
\end{remark}

\begin{example}\label{Ex:2-1}
Let ${\it Sp}(T) \cap (-\infty, 0] = \emptyset$. Then 
\begin{equation}
T^r= \frac{\sin r \pi}{\pi} \int_0^{\infty} 
 \frac{1}{\lambda} (\lambda I + T)^{-1}T\lambda ^r d\lambda \quad (0<r<1). \label{eq:4}
\end{equation}
$(T^{1/2})^2 = T$, but, needles to say, there are plenty of $X$ such that $X^2=T$. 
\end{example}
% = \frac{\sin r \pi}{\pi} \int_0^{\infty} 
% (\frac{1}{\lambda} I - (\lambda I + T)^{-1}) \lambda ^{r} d\lambda  \notag \notag 
%%%%%%%%%%%%%%%%%%%%%%%%%%%%%%%%%
\section{Accretive Operators}
Let $A$ be an accretive operator. Then 
$${\it Sp}(A)\subseteq \{z: \Re(z)\geqq 0\}, \quad  
||(A+ \lambda I)\mathbf{x}||\geq \lambda ||\mathbf{x}|| \,\,\,\, (\forall\lambda > 0, \,\, \forall\mathbf{x}).$$
 This implies that
\begin{equation}
 ||(A + \lambda I)^{-1}||\leqq \frac{1}{\lambda}. 
\label{eq:5}
\end{equation} and that   
$$\frac{1}{\lambda}I - (A + \lambda I)^{-1}$$
 is accretive too. \par
It is obvious that 
$$( A {\bf x}, {\bf x}) = 0\, \Longrightarrow \, \Re(A){\bf x} = 0\,  \Longrightarrow \,A{\bf x}= - A^*{\bf x},
$$ and  
$$\mathcal{N}(A + i \mu I) = \mathcal{N}( A^* - i\mu I)$$
for every real number $\mu$, where  
$\mathcal{N}(T):= \{\mathbf{x}: T \mathbf{x} = 0\}$. 
%It is evident that this eigenspace reduces $A$, and that if $A$ has purely imaginary %eigenvalues more, 
%then  the eigenspaces correspondent them are mutually orthogonal. 
%Denote the orthogonal sum of all these eigenspaces by $\mathcal{L}$. It reduces 
%$A$ and 
%We have a bit more general decomposition: 
%$$\mathcal{L}:= \bigcap_{n=1}^{\infty} \{ {\bf x}: A^n {\bf x} =
% (- A^*)^n {\bf x} \}.$$
%It is obvious that 
%$\mathcal{L} \supseteqq  \mathcal{N}$, $\mathcal{L}$ 
% reduces $A$ and the real part of $A|_{\mathcal{L}}$ vanishes.
%we hence get the orthogonal decomposition of $A$:  
%$$A = i H \oplus A',$$
% where $H$ is Hermitian, and $A'$ is accretive and does not have  any purely %imaginary
% eigenvalue. But $\Re(A')$ is not necessarily invertible. \\
% Let $A$ and $B$ be accretive operators. Then $A+B$ is not necessarily 
%invertible even if both of $A$ and $B$ are invertible; for instance, 
%$$ A= \left( \begin{array} {l r} 0 & i \\ i & 1\end{array}\right), \quad B= A^*.$$
%However, if $\Re(A)$ is positive definite,  $A+B$ is invertible as well.  
\begin{lemma}\label{Lem:3-1}
Let $A$ be an accretive operator and $n$ a natural number. Then 
$$\mathcal{N}(A^n) = \mathcal{N}(A).$$
\end{lemma}
\begin{proof}  $\mathcal{N}(A) \subseteqq \mathcal{N}(A^n)$ is apparent. 
For any ${\bf x} \in \mathcal{N}(A^n)$ we have $A A^{n-1} {\bf x} = {\bf 0}$. We hence 
get $A^* A A^{n-2} {\bf x} = A^* A^{n-1} {\bf x} = {\bf 0}$, which 
$A A^{n-2} {\bf x} = {\bf 0}$ follows from. By repeating this way, we arrive at  
$A {\bf x} = {\bf 0}$.
\end{proof}
Let us define a function $f(A)$ of an accretive operator $A$ by $f(z)$ given by \eqref{eq:2}. We need to 
pay attention to that the case $0\in \it{Sp}(A)$ 
is not excluded. 


\begin{theorem}\label{Th:3-1}
Let $A$ be an accretive operator and  $f(t)$ an operator monotone function
 on $(0, \infty)$ expressed by \eqref{eq:1}. 
Then   
\begin{equation}
 f(A): = f(0)I + b A + \int_0^{\infty} (\frac{1}{\lambda} I - (\lambda I + A)^{-1})
d\mu(\lambda) \label{eq:6}
\end{equation}
 is a bounded operator. \\
Let $f(A+ sI)$ be functions defined by the Riesz-Dunford integral for $s>0$. Then 
$$f(A+ sI)  \Longrightarrow f(A) \quad (s \to +0).$$
Further, $f(A)^* = f(A^*)$, and $f(A)$ remains accretive if $f(0)\geqq 0$. 
\end{theorem}

\begin{proof}
 Since \\
$||\frac{1}{\lambda} I - (\lambda I  + A)^{-1}|| \leqq  2\frac{1}{\lambda} $ and $\lim_{\lambda \to + \infty} \lambda^2 ||\frac{1}{\lambda} I - (\lambda I  + A)^{-1}|| = || A||$, 
 $$\int_0^{\infty} ||\frac{1}{\lambda} I - (\lambda I  + A)^{-1}|| d\mu(\lambda)
<\infty.$$
$f(A)$, the right side of \eqref{eq:6}, is therefore well-defined and bounded.  
Express $f(A + sI)$ as \eqref{eq:3}. Then 
\begin{align*}
&( f(A + s I) {\bf x}, {\bf y}) -( f(A) {\bf x}, {\bf y}) = 
 b  s ({\bf x}, {\bf y}) \\ 
&+ \int_0^{\infty} \left(( (\lambda  I + A)^{-1} {\bf x}, {\bf y}) - 
( ((\lambda + s) I + A)^{-1} {\bf x}, {\bf y}) \right)
 d\mu(\lambda),
\end{align*}
and consequently  
$$|| f(A + sI) - f(A)|| \leqq b  s +  \int_0^{\infty} ||(\lambda  I + A)^{-1}   - 
((\lambda + s) I + A)^{-1} || d\mu (\lambda).
$$
Write the integrand $h_s (\lambda)$. Then 
$$ h_s(\lambda) \leqq  g(\lambda) \quad  (0<s<1, \, 0<\lambda <\infty),$$
where 
$$ g(\lambda): =
\begin{cases}
\frac{2}{\lambda}  & \text{on} \,\, 0 <\lambda \leqq \frac{1}{2}\\
\frac{1}{\lambda^2}  & \text{on} \,\, \frac{1}{2}\leqq \lambda <\infty.
\end{cases}
$$
$g(\lambda)$ is obviously integrable. 
Since $\mu(\{0\}) =0$ and 
\begin{center}
$||(\lambda  I + A)^{-1}   - 
((\lambda + s) I + A)^{-1}|| \to 0$ as $s\to + 0$ for $\lambda >0$,
\end{center}
by the Lebesgue theorem,  
$$|| f(A + sI) - f(A)|| \to 0 \quad (s\to +0).$$
 From \eqref{eq:6} it directly follows that $f(A^{*})= f(A)^{*}$ and that 
$f(A)$ is accretive, because $f(0)\geqq 0$ and 
$\frac{1}{\lambda}I - (A + \lambda I)^{-1}$ is accretive.
\end{proof}

This theorem shows that \eqref{eq:4} holds for an accretive operator $A$ 
even if $0\in \it{Sp}(A)$, which is known as the formula by Balakrishnan\cite{Bal} (also see 
\cite{Kato-book, N-F}).\par 
\begin{lemma}\label{Le:2}
Let $\{A_n\}$ be a sequence of accretive operators which converges to an operator 
$A$ in the operator norm. Then for an operator monotone function 
$f(t) >0$ on $(0, \infty)$
$$f (A_n) \Longrightarrow f(A).$$
\end{lemma}

\begin{thebibliography}{99}
\bibitem[B-M]{B-M} R. Bhatia, {\it Matrix analysis}, Springer, (1997).
\bibitem[Bal]{Bal} A. V. Balakrishnan, Fractional powers of closed operators and the semigroups generated by them, Paciphic J. Math. 10 (1960) 419--437.
\bibitem[D]{D} W.\;F.\;Donoghue, {\it Monotone Matrix Functions and Analytic Continuation}. Springer-Verlag, 1974.
\bibitem[D-S]{D-S} N. Dunford, J. Schwartz, {\it Linear operators: general theory}, Wiley, New York, 1957.
\bibitem[Drury]{Drury} S. Drury, Principal powers of matrices with positive definite real part, Linear and Multilinear Alg. 63(2015) 296--301.
\bibitem[Hansen]{Hansen} F. Hansen, The fast track to L\H{o}wner's theorem, Linear Alg. Appl. 438,4557--4571(2013).
\bibitem[Kato-book]{Kato-book} T. Kato, Perturbation Theory for Linear Operators, Springer-Verlag, 1980. %
\bibitem[N-F]{N-F} B. Sz.-Nagy, C. Foias, H. Bercovici, L. Kerchy, \it Harmonic Analysis of Operators on Hilbert Space, second edition,\rm North-Holland, (1970), 2009.
\bibitem[Rudin]{Rudin} W, Rudin, \it Functional Analysis, \rm McGRW-Hill, 1974.
\bibitem[Sano]{Sano} T. Sano and K. Sagawa, Operator means for operators with positive definite real part, Adv. Operator Theory.
\bibitem[Simon]{Simon} B. Simon, Loewner's Theorem on Monotone Matrix Functions, GMW 354, Springer 2019.
\bibitem[Y]{Y} K. Yosida, \it Functional Analysis, \rm Springer, Berlin, 1965
\end{thebibliography}
\end{document}
